\section*{PHẦN KẾT LUẬN}
\addcontentsline{toc}{section}{PHẦN KẾT LUẬN}

Chặng đường hơn tám thập kỷ hình thành và phát triển của hệ thống giáo dục nghề nghiệp Việt Nam từ năm 1945 đến nay là một bản hùng ca gắn liền với những trang sử vẻ vang của dân tộc. Khởi đầu từ những lớp dạy nghề cơ khí quân giới thô sơ, cơ động giữa đại ngàn chiến khu Việt Bắc trong khói lửa kháng chiến chống thực dân Pháp, hệ thống đào tạo nghề đã không ngừng lớn mạnh, trở thành đội quân xung kích đào tạo hàng chục triệu thợ kỹ thuật lành nghề phục vụ công cuộc xây dựng chủ nghĩa xã hội ở miền Bắc, chi viện cho chiến trường miền Nam, tái thiết đất nước sau ngày non sông thống nhất, và ngày nay đang là lực lượng trụ cột vận hành nền kinh tế trong thời kỳ công nghiệp hóa, hiện đại hóa và hội nhập quốc tế.

Nghiên cứu tiến trình lịch sử qua sáu giai đoạn phát triển cho thấy giáo dục nghề nghiệp Việt Nam luôn vận động theo những quy luật khách quan sâu sắc. Mỗi bước chuyển biến về thể chế quản lý nhà nước, mỗi sự điều chỉnh về cơ cấu trình độ đào tạo đều phản ánh trực tiếp yêu cầu thích ứng với sự thay đổi của mô hình kinh tế và trật tự xã hội. Những thành tựu vĩ đại mà ngành nghề đạt được là minh chứng cho đường lối đúng đắn của Đảng và Nhà nước trong việc coi trọng giai cấp công nhân và người lao động trực tiếp; đồng thời, những điểm nghẽn cố hữu và hạn chế kéo dài qua các thời kỳ cũng để lại những bài học kinh nghiệm lịch sử vô giá về sự cần thiết phải duy trì một thể chế quản lý thống nhất, ổn định, bài học về quy hoạch cân đối tháp nhân lực quốc gia, bài học về sự gắn kết máu thịt giữa nhà trường với doanh nghiệp và bài học về công tác phân luồng học sinh phổ thông.

Bước vào kỷ nguyên mới, đứng trước những cơ hội và thách thức chưa từng có của cuộc Cách mạng công nghiệp lần thứ tư, sự phát triển bùng nổ của trí tuệ nhân tạo và xu thế chuyển dịch năng lượng xanh, giáo dục nghề nghiệp Việt Nam đang đứng trước thời khắc bản lề mang tính lịch sử. Bằng việc thấm nhuần và vận dụng sáng tạo các bài học lịch sử đã được kiểm chứng bằng thực tiễn xương máu của quá khứ, cùng với quyết tâm chính trị mạnh mẽ theo tinh thần Chỉ thị số 21-CT/TW của Ban Bí thư và Chiến lược phát triển giáo dục nghề nghiệp giai đoạn đến năm 2030, tầm nhìn đến năm 2045, hệ thống giáo dục nghề nghiệp Việt Nam hoàn toàn có đủ thế và lực để vươn mình bứt phá, trở thành bệ phóng vững chắc nâng tầm kỹ năng lao động quốc gia, góp phần hiện thực hóa khát vọng xây dựng một nước Việt Nam hùng cường, thịnh vượng và hạnh phúc.
